\subsection{李群的李代数}

设 \(\mathbf{G}\) 为李群。在 \(\mathrm{U}(\mathbf{G})\) 中，如同在任何结合代数中一样，我们记 \([t, t'] = t * t' - t' * t\)。由于 \(\mathrm{T}_{e}(\mathbf{G})\) 是 \(\mathrm{U}(\mathbf{G})\) 的本原元素之集，有 \([\mathrm{T}_{e}(\mathbf{G}), \mathrm{T}_{e}(\mathbf{G})] \subset \mathrm{T}_{e}(\mathbf{G})\) (Chapter II, § 1, no. 2, Proposition 4)。因此，括号在 \(\mathrm{T}_{e}(\mathbf{G})\) 上的限制在 \(\mathrm{T}_{e}(\mathbf{G})\) 上定义了一个李代数结构。

引理 1. 设 \(\mathbf{X}\) 与 \(\mathbf{X}'\) 为完备可赋范空间，\(\mathbf{X}_0\) 为 \(\mathbf{X}\) 中 0 的开邻域，\(f\) 为 \(\mathbf{X}_0\) 到 \(\mathbf{X}'\) 的解析映射且满足 \(f(0) = 0\)。设 \(f = f_1 + f_2 + f_3 + \dots\) 为 \(f\) 在 0 处的幂级数展开，其中 \(f_i\) 是 \(i\) 次齐次连续多项式，定义在 \(\mathbf{X}\) 上、取值于 \(\mathbf{X}'\)。设 \(t\) 为 \(\mathrm{TS}^2 (\mathbf{X})\) 的元素，把它看作 \(\mathbf{X}\) 上支集含于 \(\{0\}\) 的点分布。设 \(t' = f_{\star}(t) \in \mathrm{TS}(\mathbf{X}')\)。\(t'\) 的 1 次齐次分量是 \(\langle f_2, t \rangle\)。

设 \(t_{1}^{\prime}\) 为这个分量。于是，对 \(X^{\prime}\) 到一个多范空间中的每个连续线性映射 u，

\[
\begin{array}{l l} u (t _ {1} ^ {\prime}) = \langle t ^ {\prime}, u \rangle & \text { because   } u \text {   is   continuous   and   linear } \\ = \langle t, u \circ f \rangle & (\text { Diff.   \&   Anal.   Man.,   R,   13.2.3 }) \\ = \langle t, u \circ f _ {2} \rangle & \text { because   } t \in \mathsf {T S} ^ {2} (\mathbf {X}) \\ = u (\langle t, f _ {2} \rangle) & (\text { Diff.   \&   Anal.   Man.,   R,   13.2.2 }), \end{array}
\]

引理由此得证。

命题 24. 设 G 为李群，(U, \(\phi\) , E) 为 G 上满足 \(\phi(e) = 0\) 的卡。设 \(e\) 的开邻域 V 满足 \(\mathrm{V}^2 \subset \mathrm{U}\)。设 \(m\) 为解析映射 \((a, b) \mapsto \phi(\phi^{-1}(a)\phi^{-1}(b))\)（从 \(\phi(V) \times \phi(V)\) 到 E 中）。设

\[
m = \sum_ {i, j \geqslant 0} m _ {i, j}
\]

为 \(m\) 在 \((0,0)\) 处的幂级数展开，其中 \(m_{i,j}\) 是双次数 \((i,j)\) 的、\(\mathbf{E} \times \mathbf{E}\) 上取值于 \(\mathbf{E}\) 的双齐次连续多项式。

\begin{enumerate}
\def\labelenumi{(\roman{enumi})}
\item
  \(m_{i,0} = m_{0,j} = 0\)（对所有 \(i \neq 1\) 与 \(j \neq 1\)）。
\item
  \(m_{1,0}(a,b)=a\) 且 \(m_{0,1}(a,b)=b\)（对所有 \(a\in E\)、\(b\in E\)）。
\item
  设 \(\psi: \mathrm{T}_e(\mathrm{G}) \to \mathrm{E}\) 为 \(\phi\) 在 \(e\) 处的微分。对所有 \(u, v\)（\(\mathrm{T}_e(\mathrm{G})\) 中的元素），
\end{enumerate}

\[
\psi ([ u, v ]) = m _ {1, 1} (\psi (u), \psi (v)) - m _ {1, 1} (\psi (v), \psi (u)).
\]

\(m(a,0) = a\)，\(m(0,b) = b\)（对所有 \(a\)、\(b\)，\(\phi (\mathrm{V})\) 中的元素），这就证明了 (i) 与 (ii)。设 \(u, v\) 属于 \(\mathrm{T_e(G)}\)。把 \(\mathrm{T_0(E)}\) 与 E 认同，从而把 \(\psi\) 与 \(\mathrm{T_e}(\phi)\) 认同。\(u\) 与 \(v\) 在 \(\mathrm{T_e}(\phi)\) 下的像分别是 \(\psi (u)\) 与 \(\psi (v)\)。这些像的张量积点分布是 \((\psi(u),0)\) 与 \((0,\psi (v))\) 在 \(\mathsf{TS(E\times E)} = \mathsf{TS(E)}\times \mathsf{TS(E)}\) 中的对称积，也就是

\[
(\psi (u), 0) \otimes (0, \psi (v)) + (0, \psi (v)) \otimes (\psi (u), 0).
\]

因而 \(\phi_{*} (u * v)\) 是上述元素在映射 \(m\)（从 \(\phi(V) \times \phi(V)\) 到 \(E\) 中）下的像。它在 TS(E) 中的 1 次分量由引理 1 是

\[
x = \langle m _ {1, 1}, (\psi (u), 0) \otimes (0, \psi (v)) + (0, \psi (v)) \otimes (\psi (u), 0) \rangle .
\]

我们用下式定义一个双线性映射 \(n\colon (\mathbf{E}\times \mathbf{E})^2\to \mathbf{E}\)

\[
n ((a, b), (a ^ {\prime}, b ^ {\prime})) = m _ {1, 1} (a, b ^ {\prime}).
\]

于是 \(n((a, b), (a, b)) = m_{1,1}(a, b)\)，因而

\[
\begin{array}{l} x = \langle n, (\psi (u), 0) \otimes (0, \psi (v)) + (0, \psi (v)) \otimes (\psi (u), 0) \rangle \\ \quad = m _ {1, 1} (\psi (u), \psi (v)) + m _ {1, 1} (0, 0) = m _ {1, 1} (\psi (u), \psi (v)). \end{array}
\]

类似地，\(\phi_{*} (v * u)\) 以 \(m_{1,1}(\phi(v), \phi(u))\) 为其在 TS(E) 中的 1 次分量。由于 \(\phi([u, v])\) 是 1 次的，这就证明了 (iii)。

推论。可赋范空间 \(\mathrm{T}_e(\mathrm{G})\) 连同括号一起是一个可赋范李代数。

定义 6. 可赋范空间 \(\mathrm{T_e(G)}\) 连同括号一起称为 \(\mathbf{G}\) 的可赋范李代数，或简称为 \(\mathbf{G}\) 的李代数，记作 \(\mathbf{L(G)}\)。

命题 25. 设 \(\mathbf{G}\) 为李群，\(\operatorname{E}(\mathbf{G})\) 为 \(\mathbf{L}(\mathbf{G})\) 的包络代数。\(\mathbf{L}(\mathbf{G})\) 到 \(\operatorname{E}(\mathbf{G})\) 的典范单射定义了一个同态 \(\theta\)（从代数 \(\operatorname{E}(\mathbf{G})\) 到代数 \(\operatorname{U}(\mathbf{G})\) 中）。若 \(\mathbf{K}\) 的特征为 \(0\)，则 \(\eta\) 是双代数同构。

双代数 U(G) 是余交换的 (Differentiable and Analytic Manifolds, R, 13.5.1)，且滤过 (U \(_{s}\) (G)) 与双代数结构相容。U(G) 的本原元素之集是 L(G)。于是只需引用 Chapter II, § 1, no. 6, Theorem 1。

当 \(\mathbf{K}\) 的特征为 0 时，我们今后将把 \(\mathrm{U}(\mathrm{G})\) 与 L(G) 的包络代数认同。由 (2) 与命题 7 (ii)，U(G) 到 U(G) 的映射 \(t \mapsto t^{\vee}\) 这时与 U(G) 的主反自同构认同 (Chapter I, § 2, no. 4)。

命题 26. 设 K 的特征为 \(p > 0\)。对所有 \(a \in \mathrm{L}(\mathrm{G})\)，有 \(a^p \in \mathrm{L}(\mathrm{G})\) 且 \(\operatorname{ad}(a^p) = (\operatorname{ad} a)^p\)（幂 \(a^p\) 是在 \(\mathrm{U}(\mathrm{G})\) 中计算的）。

若 \(a \in \mathrm{L}(\mathrm{G})\)，则 \(a\) 在 \(\mathrm{U}(\mathrm{G})\) 中是本原的，因而 \(a^p\) 在 \(\mathrm{U}(\mathrm{G})\) 中是本原的 (Chapter II, § 1, no. 2, Remark 1)，于是 \(a^p \in \mathrm{L}(\mathrm{G})\)。设 \(\sigma_a\)（相应地 \(\tau_a\)）为线性映射 \(x \mapsto a * x\)（相应地 \(x \mapsto x * a\)），它从 \(\mathrm{U}(\mathrm{G})\) 到 \(\mathrm{U}(\mathrm{G})\) 中。对所有 \(x \in \mathrm{U}(\mathrm{G})\)，(ad \(a\) ) \((x) = (\sigma_a - \tau_a)(x)\)，因而 (ad \(a\) ) \(^p = (\sigma_a - \tau_a)^p\)。但 \(\sigma_a\) 与 \(\tau_a\) 交换，故 \((\sigma_a - \tau_a)^p = (\sigma_a)^p - (\tau_a)^p = \tau_a^p - \sigma_a^p\)，由此得第二个断言。

定义 7. 设 X 为 \(\mathbf{C}^r\) 类流形（\(r \geqslant 2\)），\(\mathfrak{g}\) 为完备可赋范李代数。一个 \(\mathbf{C}^{r-1}\) 类的、\(\mathfrak{g}\) 在 X 上的无穷小左（相应地右）作用律是指映射 \(a \mapsto \mathbf{D}_a\)（从 \(\mathfrak{g}\) 到 X 上向量场的集合），它满足下列性质：

\begin{enumerate}
\def\labelenumi{(\alph{enumi})}
\item
  映射 \((a, x) \mapsto \mathrm{D}_a(x)\) 是一个 \(\mathbf{C}^{r-1}\) 类态射（从平凡向量丛 \(g \times X\) 到向量丛 \(\mathrm{T}(X)\)）；
\item
  \([\mathrm{D}_a,\mathrm{D}_b] = -\mathrm{D}_{[a,b]}\)（相应地 \([\mathrm{D}_a,\mathrm{D}_b] = \mathrm{D}_{[a,b]})\)）（对所有 \(a,b\)，\(\mathfrak{g}\) 中的元素）。
\end{enumerate}

特别地，每个向量场 \(\mathbf{D}_a\) 都是 \(\mathbf{C}^{r - 1}\) 类的。

注。设 X 为 \(C^{r}\) 类流形，g 为有限维李代数，\(a \mapsto D_{a}\) 为 g 到 X 上 \(C^{r-1}\) 类向量场的向量空间中的一个线性映射。于是定义 7 的条件 (a) 成立。因为，通过取 g 的一个基并引用 Differentiable and Analytic Manifolds, R, 7.7.1，问题就化为 dim g = 1 的情形，而这时我们的断言是明显的。

命题 27. 设 G 为李群，X 为 \(\mathbf{C}^r\) 类流形。设给定了 G 在 X 上的一个 \(\mathbf{C}^r\) 类左（相应地右）作用律。对所有 \(a \in \mathrm{L}(\mathrm{G})\)，用 \(\mathrm{D}_a\) 表示由 \(a\) 在 X 上定义的点分布场。

\begin{enumerate}
\def\labelenumi{(\roman{enumi})}
\item
  映射 \((a, x) \mapsto \mathrm{D}_a(x)\) 是一个 \(\mathbf{C}^{r-1}\) 类态射（从平凡向量丛 \(\mathbf{L}(\mathbf{G}) \times \mathbf{X}\) 到向量丛 \(\mathrm{T}(\mathbf{X})\)）。
\item
  设 \(\mathbf{I}\) 为 \(\mathbf{K}\) 中包含 0 的开子集，\(\gamma: \mathbf{I} \to \mathbf{G}\) 为 \(\mathbf{C}^r\) 类映射且满足 \(\gamma(0) = e\)。设 \(a = \mathrm{T}_0(\gamma) 1 \in \mathbf{L}(\mathbf{G})\)。若 \(f\) 是 \(\mathbf{C}^r\) 类函数，定义在 \(\mathbf{X}\) 的某个开子集上，则
\end{enumerate}

\((D_{a}f)(x) = \lim_{k\in \mathbb{K}^{*},k\to 0}k^{-1}(f(\gamma (k)x) - f(x))\)（若 G 左作用），

\[
(D _ {a} f) (x) = \lim _ {k \in K ^ {*}, k \rightarrow 0} k ^ {- 1} (f (x \gamma (k)) - f (x)) \quad \text {if } G \text { operates on the right}.
\]

\begin{enumerate}
\def\labelenumi{(\roman{enumi})}
\setcounter{enumi}{2}
\tightlist
\item
  若 \(r \geqslant 2\)，则映射 \(a \mapsto \mathbf{D}_a\) 是一个 \(\mathbf{C}^{r-1}\) 类的、\(\mathbf{L}(\mathbf{G})\) 在 \(\mathbf{X}\) 上的无穷小左（相应地右）作用律。
\end{enumerate}

设 G 左作用在 X 上。设 \(\phi: G \times X \to X\) 为该作用律。于是 \(T(\phi)\) 是一个 \(\phi\)-态射，为 \(C^{r-1}\) 类，从向量丛 \(T(G) \times T(X)\) 到向量丛 \(T(X)\) 中 (Differentiable and Analytic

Manifolds, R, 8.1.2)。诱导向量丛 \((\mathrm{T}(\mathrm{G}) \times \mathrm{T}(\mathrm{X}))|(\{e\} \times \mathrm{X})\) 与 \(\mathrm{E} = \mathrm{L}(\mathrm{G}) \times \mathrm{T}(\mathrm{X})\) 认同。因而 \(\mathrm{T}(\phi)|\mathrm{E}\) 是 \(\mathbf{C}^{r-1}\) 类向量丛态射。对 \((a, x) \in \mathrm{L}(\mathrm{G}) \times \mathrm{X}\)，有 \(\mathrm{T}(\phi)(a, x) = \mathbf{D}_{a}(x)\)，由此得 (i)。

给出 \((\mathrm{D}_a f)(x)\) 的公式由 §2 第 2 节末尾以及 Differentiable and Analytic Manifolds, R, 8.4.5 得出。

设 \(r \geqslant 2\)。设 \(a, b\) 属于 \(\mathbf{L}(\mathbf{G})\)，\(f\) 为 \(\mathbf{C}^r\) 类函数，定义在 \(\mathbf{X}\) 的某个开子集上。于是

\[
\begin{array}{l l} \mathrm{D} _ {[ a, b ]} f = \mathrm{D} _ {b} (\mathrm{D} _ {a} f) - \mathrm{D} _ {a} (\mathrm{D} _ {b} f) & \text { by   (17) } \\ = [ \mathrm{D} _ {b}, \mathrm{D} _ {a} ] f & (\text { Diff.   \&   Anal.   Man.,   R,   8.5.3 }). \end{array}
\]

设 \(x \in X\)。取 \(f\) 为 \(x\) 的某开邻域上的一个卡，便得 \(D_{[a,b]}(x) = [D_b, D_a](x)\)，由此得 (iii)。若 \(G\) 右作用在 \(X\) 上，论证是类似的。

当 \(r \geqslant 2\) 时，映射 \(a \mapsto D_{a}\) 称为与所给作用律相伴的无穷小作用律。

